\documentclass[11pt,a4paper]{article}

\usepackage[T1]{fontenc}
\usepackage[utf8]{inputenc}
\usepackage{lmodern}
\usepackage[french]{babel}
\usepackage{amsmath,amssymb,mathtools}
\usepackage[margin=2.4cm]{geometry}
\usepackage{booktabs,array,enumitem,xcolor,parskip,setspace}
\usepackage[colorlinks=true,linkcolor=blue!50!black,urlcolor=blue!40!black]{hyperref}
\usepackage{algorithm}
\usepackage[noend]{algpseudocode}
\usepackage{tcolorbox}
\tcbuselibrary{skins,breakable}

\setstretch{1.08}
\emergencystretch=2em
\setlist{itemsep=2pt,topsep=3pt}

\newcommand{\fic}[1]{\texttt{\detokenize{#1}}}
\newtcolorbox{enclair}[1][]{colback=blue!4!white, colframe=blue!35!black, boxrule=0.5pt, arc=2pt,
  breakable, title={#1}, fonttitle=\bfseries, left=6pt, right=6pt, top=4pt, bottom=4pt}

\newcommand{\E}{\mathbb{E}}
\newcommand{\Mt}{\widetilde{M}}
\newcommand{\hb}{\bar h}
\newcommand{\mh}{\hat m}
\DeclareMathOperator{\sign}{sign}
\DeclareMathOperator{\rms}{rms}
\DeclareMathOperator{\eigvecs}{vecp}
\DeclareMathOperator{\med}{med}

\floatname{algorithm}{Algorithme}
\renewcommand{\algorithmicrequire}{\textbf{Entrées :}}
\renewcommand{\algorithmicensure}{\textbf{Sortie :}}
\algrenewcommand\algorithmicif{\textbf{si}}
\algrenewcommand\algorithmicthen{\textbf{alors}}
\algrenewcommand\algorithmicelse{\textbf{sinon}}
\algrenewcommand\algorithmicfor{\textbf{pour}}
\algrenewcommand\algorithmicdo{\textbf{faire}}

\title{\Large Écrêter au lieu d'amortir :\\ l'expression du clipping et l'algorithme d'entraînement\\[6pt]
\large Famille B de l'expérience E16 (\fic{docs/reports/plan_floor_clip.md})}
\author{}
\date{Septembre 2026}

\begin{document}
\maketitle

\begin{enclair}[En bref]
Les modes \fic{ekfac} et \fic{tekfac} divisent chaque coordonnée du pas, dans une base propre, par
une courbure stockée à laquelle on ajoute une constante de sécurité $\lambda$. La famille B remplace
cette addition par un \emph{écrêtage} : on divise par la courbure seule, puis on plafonne chaque
coordonnée à un seuil $\gamma$. Le seuil est toujours fixé \emph{couche par couche}, parce que
l'erreur d'échelle de la courbure stockée varie d'une couche à l'autre jusqu'à quatre ordres de
grandeur. Trois façons de le fixer sont implémentées, parce qu'elles ne posent pas la même question :
\begin{itemize}
  \item \textbf{quantile} : $\gamma$ recalculé à chaque pas pour écrêter une proportion fixe $q$ des
        coordonnées. C'est une normalisation du pas par couche autant qu'un écrêtage ;
  \item \textbf{ema} (bras principal) : le même seuil, corrigé par la taille du momentum rapportée à sa
        moyenne (corrigée du biais) sur environ $1\,000$ pas, pour que le pas suive de nouveau le
        momentum ;
  \item \textbf{fixe} : la médiane du seuil de la couche sur les $1\,000$ pas précédant le pas
        $2\,000$, puis gelée. C'est un écrêtage conditionnel, comme celui de Sophia.
\end{itemize}
C'est l'option \fic{rescale_form="clip"} d'\fic{AdaFisherMulti}, avec \fic{clip_threshold} et
\fic{clip_fraction=q}, implémentée dans\\ \fic{src/adafisher_modes/approximations/_rescale_utils.py}.
\end{enclair}

\section{Notations}

Tout est écrit pour \fic{ekfac}. Pour \fic{tekfac}, il suffit de remplacer $Q_A \to Q_\Phi$,
$Q_B \to Q_\Psi$ et $s^\star \to \Theta$.

Pour une couche $\ell$, on appelle \emph{ligne} un couple (exemple, position) : une seule position
pour une couche dense appliquée à un vecteur, $T_\ell$ positions pour une convolution (les positions
de sortie) ou pour une couche appliquée à une séquence de $T_\ell$ jetons. Pour chaque ligne $t$ :
\begin{itemize}
  \item $\hb_{t}$ est l'entrée de la couche, sous la forme que lui donne l'optimiseur :
        \begin{itemize}
          \item couche dense ou convolution \emph{avec} biais : l'entrée (le patch, pour une
                convolution), augmentée d'un $1$ pour le biais ;
          \item couche \emph{sans} biais : l'entrée seule, sans le $1$ ;
          \item couche de normalisation (\fic{LayerNorm}, \fic{BatchNorm2d}) :
                $\hb_t = (z_t,\ 1)$, où $z_t$ est la moyenne sur les canaux de l'entrée \emph{brute},
                pas de l'activation normalisée (déviation 4.6 de \fic{CLAUDE.md}). Les paramètres
                d'échelle et de décalage forment alors une matrice $C \times 2$ ;
        \end{itemize}
  \item pour une couche de normalisation, $\hat x_t \in \mathbb R^C$ est l'activation normalisée de la
        ligne, recalculée avec les statistiques qu'utilise la couche : par jeton pour
        \fic{LayerNorm} ; pour \fic{BatchNorm2d}, celles du batch en entraînement et les statistiques
        courantes en évaluation. Le vrai gradient de la ligne par rapport à (échelle, décalage) est
        la matrice $C\times 2$ $\,[\delta_t \odot \hat x_t,\ \delta_t]$ ;
  \item $\delta_{t}$ est le gradient de la perte moyenne du batch par rapport à la sortie de la couche
        à cette ligne ;
  \item $A = \E[\hb\hb^\top]$ et $B = \E[\delta\delta^\top]$, moyennes sur les $N\,T_\ell$ lignes du
        batch, sont les deux facteurs de Kronecker, tenus en moyenne glissante ; $Q_A$ et $Q_B$ sont
        leurs vecteurs propres ;
  \item $s^\star$, de la forme de la matrice des paramètres ($d_\text{out} \times d_\text{in}$, avec
        une colonne de plus pour le biais s'il existe), est la courbure stockée dans la base propre
        $Q_B \otimes Q_A$ (le « rééchelonnement » d'EKFAC) ;
  \item $m$ est le momentum de la couche, poids et biais réunis dans une matrice de cette même forme,
        $n$ son nombre de coordonnées, et $\mh_t = m_t/(1-\beta^{t+1})$ le momentum corrigé du biais.
\end{itemize}

\textbf{Une conséquence des lignes : l'échelle de $s^\star$ dépend de la couche.} Moyenner sur
$N\,T_\ell$ lignes au lieu de $N$ exemples divise la courbure stockée d'une couche partagée par un
facteur de l'ordre de $T_\ell$ (\fic{plan_exp_lot3.md} §5.2). Le surrogate des couches de
normalisation ajoute une erreur de plus. Mesuré au pas $2\,000$ avec ce surrogate, par rapport à la
tête : de $1/422$ à $1/7\,400$ sur les convolutions de \fic{cnn_gn_cifar}, de $1/251$ à $1/5\,620$ sur
\fic{vit_micro_cifar}, de $1/91$ à $1/2\,724$ sur \fic{cct_2_3x2_cifar}. C'est pourquoi aucun seuil
n'est partagé entre couches. L'option \fic{norm_exact_rescaling}, activée dans tous les bras d'E16,
supprime la seconde erreur (§\ref{sec:remarques}), pas la première.

On projette le momentum corrigé dans la base propre :
\[
  \Mt = Q_B^\top\, \mh\, Q_A .
\]
Toutes les opérations qui suivent sur $\Mt$ et $s^\star$ sont faites coordonnée par coordonnée.

\section{Trois façons de diviser par la courbure}

Dans chaque cas, la couche se déplace de $-\eta\, Q_B\, U\, Q_A^\top$, et seule la matrice $U$
change :

\begin{center}
\renewcommand{\arraystretch}{1.4}
\begin{tabular}{@{}lll@{}}
\toprule
forme & $U_i$ & taux d'apprentissage \\
\midrule
\fic{add} (correctif E14, défaut) & $\dfrac{\Mt_i}{s^\star_i + \lambda}$ & $\eta = \text{cap}\cdot\lambda$ \\[6pt]
\fic{floor} (famille A) & $\dfrac{\Mt_i}{\max(s^\star_i,\ \lambda)}$ & $\eta = \text{cap}\cdot\lambda$ \\[6pt]
\fic{clip} (famille B) & voir \eqref{eq:U} & $\eta$ du benchmark \\
\bottomrule
\end{tabular}
\end{center}

Pour \fic{add} et \fic{floor}, $\Mt$ est projeté depuis le momentum brut $m$, et la correction de
biais est appliquée ensuite au pas, comme dans AdaFisher ; l'opérateur étant linéaire, cela revient
au même. Dans ces deux formes, $\lambda$ fixe à la fois le plancher de la courbure et, avec
$\eta = \text{cap}\cdot\lambda$, la taille maximale du pas. Dans la troisième, $\lambda$ n'intervient
plus du tout.

\section{L'expression du clipping}

\subsection{L'opération d'écrêtage}

Pour une couche, à un pas donné :
\begin{align}
  r_i &= \frac{|\Mt_i|}{s^\star_i}
      && \text{taille du pas \emph{non amorti} de la coordonnée } i, \label{eq:r}\\
  g &= \kappa\ \rms(\Mt) = \kappa\Bigl(\tfrac1n\textstyle\sum_j \Mt_j^2\Bigr)^{1/2},\quad \kappa = 10^{-3}
      && \text{garde,} \label{eq:guard}\\
  \mathcal A &= \{\, i : |\Mt_i| > g \,\}, \quad n_a = |\mathcal A|
      && \text{coordonnées actives,} \label{eq:active}\\
  U_i &=
  \begin{cases}
    \sign(\Mt_i)\,\min\Bigl(\dfrac{r_i}{\gamma},\ 1\Bigr) & i \in \mathcal A, \\[10pt]
    \dfrac{\Mt_i}{\max(\gamma\, s^\star_i,\ g)} & i \notin \mathcal A.
  \end{cases} \label{eq:U}
\end{align}
\begin{itemize}
  \item Une coordonnée active est \emph{écrêtée} quand $r_i \ge \gamma$ : $U_i = \pm1$, elle bouge
        exactement de $\eta_t$. L'écrêtage est \emph{défini} par ce test, sur le tenseur $r$ même dont
        $\gamma$ est tiré ; le nombre de coordonnées écrêtées est donc exact par construction (aux
        ex-æquo près).
  \item Une coordonnée active non écrêtée vaut $\Mt_i/(\gamma\, s^\star_i)$ : le pas de Newton non
        amorti, divisé par $\gamma$. Elle garde la forme que dicte la courbure.
  \item Une coordonnée inactive bouge d'au plus $|\Mt_i|/g < 1$.
  \item Pour une coordonnée active, $\operatorname{clip}\bigl(x/(\gamma s),\,1\bigr) = x/\max(\gamma s,\,|x|)$ :
        l'écrêtage est un plancher sur la courbure, mais un plancher proportionnel au momentum
        lui-même ($|\Mt_i|/\gamma$), et non une constante $\lambda$.
\end{itemize}

\textbf{Lien avec Sophia} (Liu et al., arXiv:2305.14342). La mise à jour de Sophia est
$\theta \leftarrow \theta - \eta\,\operatorname{clip}\bigl(m/(\rho\,h + \epsilon),\,1\bigr)$.
Le $\gamma$ ci-dessus joue le rôle de $\rho$. Sophia fixe un seul $\rho$ pour tout le réseau, ce qui a
un sens chez elle parce que son estimation $h$ a la bonne échelle dans toutes les couches (§2.3) ;
ici ce n'est pas le cas. Sophia surveille la proportion de coordonnées non écrêtées (son
\fic{win_rate}, entre $0{,}1$ et $0{,}5$, soit une proportion écrêtée de $0{,}5$ à $0{,}9$).

\subsection{La garde et les coordonnées actives}

La proportion écrêtée $q$ ne compte que les coordonnées actives. Ce choix a été imposé par une mesure.
Dans la base propre, beaucoup de coordonnées sont nulles en arithmétique exacte : le gradient vit dans
l'espace engendré par les entrées du batch, et $A$ n'a que peu de directions observées. En float32,
ces coordonnées valent du bruit d'arrondi. Mesuré sur un batch de 6 exemples : $|\Mt_i|$ descend
jusqu'à $10^{-13}$, et la médiane de $r$ est entièrement fixée par ces coordonnées. Si on les
comptait, le seuil tomberait dans le bruit et toutes les coordonnées réelles seraient écrêtées. Sans
garde, une direction de courbure exactement nulle (les directions de la tête qui ajoutent une même
constante à tous les logits) ferait un pas entier piloté par l'arrondi.

\textbf{Le niveau $\kappa = 10^{-3}$ est un choix, pas un écart mesuré.} Sur les vraies couches, des
coordonnées réelles descendent jusqu'à $5\cdot10^{-6}\,\rms$ et des lignes nulles montent jusqu'à
$9{,}8\cdot10^{-4}\,\rms$. Tout verdict sur l'écrêtage est donc un verdict à $\kappa = 10^{-3}$.

\subsection{Trois façons de fixer le seuil $\gamma$, toutes couche par couche}

On note $\operatorname{quant}_q(r)$ la $c$-ième plus grande valeur de $\{r_i : i \in \mathcal A\}$,
avec $c = \max(\lfloor q\, n_a \rfloor, 1)$. Elle est calculée par un tri et une lecture d'indice sur
le GPU, sans que $n_a$ ne remonte à l'hôte.

\paragraph{Quantile : recalculé à chaque pas.}
\[
  \gamma_t = \operatorname{quant}_q(r_t) .
\]
Exactement $c$ coordonnées actives sont écrêtées à chaque pas. Pour tous $a, b > 0$,
$U(a\Mt,\ b\,s^\star) = U(\Mt,\ s^\star)$ : aucune erreur d'échelle de la courbure stockée n'atteint
le pas. Mais la taille du momentum non plus. La plus grande coordonnée de chaque couche bouge de
$\eta_t$ à chaque pas, même quand le gradient s'éteint. C'est donc une \emph{normalisation du pas par
couche} avec un écrêtage de sa forme. Cette règle sert de contrôle.

\paragraph{EMA : le pas suit de nouveau le momentum.}
Avec $\mu_t = \rms(\Mt_t)$ et $w = 1/H$, $H = 1\,000$ pas :
\[
  \nu_t = (1-w)\,\nu_{t-1} + w\log\mu_t, \qquad
  d_t = (1-w)\,d_{t-1} + w, \qquad
  \log\bar\mu_t = \frac{\nu_t}{d_t}, \qquad
  \gamma_t = \operatorname{quant}_q(r_t)\cdot\frac{\bar\mu_t}{\mu_t},
\]
avec $\nu_0 = d_0 = 0$ : c'est la moyenne glissante corrigée du biais d'Adam, appliquée à $\log\mu$.
Un pas de momentum nul ne porte pas de taille et n'est pas compté.
\begin{itemize}
  \item Si le momentum a sa taille habituelle ($\mu_t = \bar\mu_t$), c'est exactement la règle
        quantile.
  \item S'il est $a$ fois plus petit que sur les $\sim H$ derniers pas, les coordonnées non écrêtées
        sont $a$ fois plus petites et moins de coordonnées sont écrêtées. Un pic est écrêté plus fort.
  \item L'échelle de la courbure reste absorbée à chaque pas : seule la taille du momentum est
        retardée.
  \item La moyenne porte sur $\mu$ et non sur $\gamma$, parce que $\gamma$ contient aussi l'échelle de
        $s^\star$, qui saute à chaque mise à jour des facteurs ($s^\star$ est à 92 \% un seul
        minibatch).
  \item \textbf{Pourquoi la correction de biais.} Le premier momentum corrigé est un seul gradient,
        mesuré $4{,}4$ à $5$ fois plus grand en moyenne quadratique que le momentum une fois installé.
        Une moyenne initialisée avec lui restait décalée de plus de 10 \% pendant 12 à 18 \% d'un
        entraînement.
\end{itemize}

\paragraph{Fixe : un écrêtage conditionnel.}
\[
  \gamma_t =
  \begin{cases}
    \operatorname{quant}_q(r_t) & t < t_\text{cal}, \\[3pt]
    \med\bigl\{\operatorname{quant}_q(r_\tau) : t_\text{cal} - W \le \tau < t_\text{cal}\bigr\}
      \ \text{gelée} & t \ge t_\text{cal},
  \end{cases}
\]
avec $t_\text{cal} = 2\,000$ et $W = 1\,000$ (médiane inférieure). Après la calibration, l'écrêtage
n'agit que là où le pas non amorti dépasse le plafond, et en dessous le pas est proportionnel au
momentum.
\begin{itemize}
  \item \textbf{Une médiane sur une fenêtre}, parce que le quantile d'un seul pas est un tirage :
        il variait d'un facteur $3{,}6$ d'une mise à jour des facteurs à l'autre.
  \item \textbf{Couche par couche}, à cause de l'erreur d'échelle propre à chaque couche. Cela rend
        aussi la bascule continue : la valeur gelée est celle que la couche utilisait. Mesuré sur de
        vrais réseaux : la taille du pas change d'un facteur $0{,}92$ à $1{,}04$ à la bascule.
  \item $t_\text{cal}$ et le début de la fenêtre sont au-delà des $10\cdot T_\text{Cov} = 1\,000$ pas
        nécessaires à l'effacement de la graine identité (époque 2 sur 15 à batch 32).
  \item Après la calibration, la proportion écrêtée dérive librement ; elle est journalisée. C'est le
        diagnostic de Sophia.
\end{itemize}

\subsection{Propriétés communes}

\begin{enumerate}
  \item \textbf{Pas borné.} $|U_i| \le 1$ : chaque coordonnée de la base propre bouge d'au plus
        $\eta_t$, quelle que soit la courbure.
  \item \textbf{Momentum corrigé du biais.} L'écrêtage compare le momentum à un seuil, ce qui n'est
        pas une opération linéaire. Il reçoit donc $\mh$, et le pas est appliqué sans nouvelle
        division par $1-\beta^{t+1}$. Pour la règle quantile, c'est sans effet (invariance d'échelle).
  \item \textbf{Les deux limites.} Quand $q \to 1$, presque tout est écrêté : $U \approx \sign(\Mt)$,
        une descente de signe dans la base propre. Quand $q \to 0$, $U \approx \Mt/(\gamma s^\star)$ :
        le pas EKFAC non amorti, dont la plus grande coordonnée vaut $1$ (au moment de la
        calibration, pour la règle fixe).
\end{enumerate}

\section{L'algorithme d'entraînement}

Constantes livrées : moyenne glissante $X \leftarrow \rho_\text{old} X + \rho_\text{new} X_\text{inst}$
avec $(\rho_\text{old}, \rho_\text{new}) = (0{,}08,\ 0{,}008)$ ; $T_\text{Cov} = T_\text{eig} = 100$ ;
$\beta = 0{,}9$. Le calcul de la base propre se fait dans le hook backward, avant de projeter le
gradient (\fic{eig_before_rescale=True}), pour que $s^\star$ soit mesuré dans la base qu'utilise
ensuite le pas. Dans l'algorithme, $R_\ell = N\,T_\ell$ est le nombre de lignes de la couche $\ell$
dans un batch.

\begin{algorithm}[H]
\caption{EKFAC avec écrêtage (famille B), $T$ pas d'entraînement}
\label{alg:clip}
\begin{algorithmic}[1]
\Require paramètres $\theta$ ; taux de base $\eta$ et calendrier cosinus $\eta_t$ ; $\beta = 0{,}9$ ;
  proportion écrêtée $q \in (0,1)$ ; $\kappa = 10^{-3}$ ; règle $\in\{$quantile, ema, fixe$\}$ ;
  $H = 1\,000$ ; $t_\text{cal} = 2\,000$ ; $W = 1\,000$ ; weight decay $w_d$ (couplé ou découplé)
\State pour chaque couche hookée $\ell$ : $A_\ell \gets I$, $B_\ell \gets I$, $m_\ell \gets 0$,
       $\nu_\ell \gets 0$, $d_\ell \gets 0$, $\mathcal W_\ell \gets \varnothing$,
       $\gamma^\text{fixe}_\ell \gets$ non défini
\For{$t = 0, \dots, T-1$}
  \State tirer un batch de $N$ exemples ; forward
  \If{$t \bmod T_\text{Cov} = 0$}
    \State chaque couche $\ell$ garde ses $R_\ell$ lignes $\hb_{\ell,r}$ \Comment{entrées, forme de §1}
    \State et, sous \fic{norm_exact_rescaling}, chaque couche de normalisation ses $\hat x_{\ell,r}$
  \EndIf
  \State $L \gets \frac1N\sum_n \ell(f_\theta(x_n), y_n)$ ; backward
  \If{$t \bmod T_\text{Cov} = 0$} \Comment{hooks backward, pour chaque couche $\ell$}
    \State $A_\ell \gets \rho_\text{old}A_\ell + \rho_\text{new}\,\tfrac1{R_\ell}\sum_r \hb_{\ell,r}\hb_{\ell,r}^\top$
    \State $B_\ell \gets \rho_\text{old}B_\ell + \rho_\text{new}\,\tfrac1{R_\ell}\sum_r \delta_{\ell,r}\delta_{\ell,r}^\top$
    \If{$t \bmod T_\text{eig} = 0$}
      \State $Q_{A,\ell} \gets \eigvecs(A_\ell + \varepsilon I)$, \ $Q_{B,\ell} \gets \eigvecs(B_\ell + \varepsilon I)$
             \Comment{$\varepsilon = 10^{-6}\,\overline{\operatorname{diag}}$}
      \If{premier calcul} $s^\star_\ell \gets \mathbf 1$ \EndIf
    \EndIf
    \If{$\ell$ est une couche de normalisation \textbf{et} \fic{norm_exact_rescaling}}
      \State $G_\ell \gets \tfrac1{R_\ell}\sum_r \bigl(Q_{B,\ell}^\top\,
             [\delta_{\ell,r}\odot\hat x_{\ell,r},\ \delta_{\ell,r}]\;Q_{A,\ell}\bigr)^{\odot 2}$
             \Comment{vrai gradient de la ligne}
    \Else
      \State $G_\ell \gets \tfrac1{R_\ell}\sum_r \bigl(Q_{B,\ell}^\top \delta_{\ell,r}\bigr)^{\odot 2}\,
             \bigl(Q_{A,\ell}^\top \hb_{\ell,r}\bigr)^{\odot 2\,\top}$ \Comment{carrés coordonnée par coordonnée}
    \EndIf
    \State $s^\star_\ell \gets \rho_\text{old}\,s^\star_\ell + \rho_\text{new}\,G_\ell$
  \EndIf
  \For{chaque couche hookée $\ell$, poids et biais ensemble}
    \State $g_\ell \gets \nabla_{[W_\ell,\,b_\ell]} L$ \ ($+\,w_d\,[W_\ell, b_\ell]$ si weight decay couplé)
    \State $m_\ell \gets \beta\, m_\ell + (1-\beta)\, g_\ell$ ; \ $\mh_\ell \gets m_\ell / (1-\beta^{t+1})$
    \State $\Mt \gets Q_{B,\ell}^\top\, \mh_\ell\, Q_{A,\ell}$ ; \ $r \gets |\Mt| \oslash s^\star_\ell$ ;
           \ $g \gets \kappa\,\rms(\Mt)$ ; \ $\mathcal A \gets \{i : |\Mt_i| > g\}$
    \If{règle $=$ fixe \textbf{et} $\gamma^\text{fixe}_\ell$ défini}
      \State $\gamma \gets \gamma^\text{fixe}_\ell$
    \Else
      \State $\gamma \gets \operatorname{quant}_q\bigl(r_{\mathcal A}\bigr)$
             \Comment{$c$-ième plus grande, $c = \max(\lfloor q\,|\mathcal A|\rfloor, 1)$}
      \If{règle $=$ ema \textbf{et} $\rms(\Mt) > 0$}
        \State $\mu \gets \rms(\Mt)$ ; \ $\nu_\ell \gets (1-\tfrac1H)\nu_\ell + \tfrac1H\log\mu$ ;
               \ $d_\ell \gets (1-\tfrac1H)d_\ell + \tfrac1H$
        \State $\gamma \gets \gamma\cdot\exp(\nu_\ell/d_\ell)/\mu$
      \EndIf
      \If{règle $=$ fixe \textbf{et} $t_\text{cal} - W \le t < t_\text{cal}$} $\mathcal W_\ell \gets \mathcal W_\ell \cup \{\gamma\}$ \EndIf
      \If{règle $=$ fixe \textbf{et} $t \ge t_\text{cal}$}
        \State $\gamma^\text{fixe}_\ell \gets \med(\mathcal W_\ell)$ (ou $\gamma$ si $\mathcal W_\ell = \varnothing$) ;
               \ $\gamma \gets \gamma^\text{fixe}_\ell$ \Comment{gelé}
      \EndIf
    \EndIf
    \State $U_i \gets \sign(\Mt_i)\min(r_i/\gamma, 1)$ si $i \in \mathcal A$, sinon
           $U_i \gets \Mt_i/\max(\gamma s^\star_{\ell,i}, g)$ \Comment{éq.~\eqref{eq:U}}
    \State $D_\ell \gets Q_{B,\ell}\, U\, Q_{A,\ell}^\top$
    \If{weight decay découplé} $[W_\ell, b_\ell] \gets (1-\eta_t w_d)\,[W_\ell, b_\ell]$ \EndIf
    \State $[W_\ell, b_\ell] \gets [W_\ell, b_\ell] - \eta_t\, D_\ell$ \Comment{sans correction de biais}
  \EndFor
  \State autres paramètres : momentum simple avec correction de biais \Comment{gelés dans E16}
  \State $\eta_{t+1} \gets$ calendrier cosinus
\EndFor
\end{algorithmic}
\end{algorithm}

\section{Remarques}
\label{sec:remarques}

\begin{itemize}
  \item \textbf{Ce que fait la courbure ici.} Elle fixe la \emph{forme} du pas dans les trois règles.
        Sa taille est bornée par $\eta_t$ coordonnée par coordonnée. Avec la règle quantile, la taille
        du pas ne dépend plus que de $\eta_t$ ; avec les règles EMA et fixe, elle suit aussi le
        momentum.
  \item \textbf{Ce que chaque comparaison isole dans E16.} EMA contre quantile : le même écrêtage, avec
        ou sans normalisation à chaque pas. Fixe contre \fic{add} : un plancher proportionnel au
        momentum contre un plancher constant, l'énoncé le plus net de la famille B. Fixe contre EMA :
        un seuil gelé contre un seuil qui suit le momentum, tous deux par couche.
  \item \textbf{Couches de normalisation.} Sans option, leur $s^\star$ est estimé depuis
        $\hb_t = (z_t, 1)$, construit sur la moyenne des canaux de l'entrée brute. Ce n'est donc pas
        le second moment du vrai gradient de la ligne, $[\delta_t\odot\hat x_t,\ \delta_t]$. Mesuré au
        pas $300$, la colonne qui porte l'échelle est $295$ à $989$ fois trop petite sur les
        LayerNorm de \fic{vit_micro_cifar}, $80$ à $106$ fois sur celles de \fic{cct_2_3x2_cifar} ;
        la colonne de décalage est juste ($\times 0{,}96$ à $1{,}03$). L'écrêtage l'exposait à
        $q \le 0{,}3$, en freinant les pas de décalage d'un facteur $3$ à $117$.
        L'option \fic{norm_exact_rescaling} corrige cela : $s^\star$ est la moyenne des carrés du
        vrai gradient de chaque ligne, projeté dans la même base (algorithme~\ref{alg:clip}). Le
        lemme~1 d'EKFAC en fait la diagonale optimale dans cette base, quelle que soit la base. La
        base elle-même vient toujours de $A$, donc du surrogate. L'option est activée dans tous les
        bras d'E16, y compris le correctif E14 qui y est relancé.
  \item \textbf{Convolutions.} $m$ est le poids remis en forme
        $C_\text{out}\times C_\text{in}k_hk_w$ (plus une colonne de biais s'il existe). Avec SUA, $\Mt$
        a la forme $k_hk_w\times C_\text{out}\times C_\text{in}$ (même remarque), et $\gamma$ est
        calculé sur l'ensemble des positions du noyau de la couche.
  \item \textbf{Coût.} Par couche et par pas, en plus d'EKFAC : un tri sur $n$ valeurs, une lecture
        d'indice et quelques opérations coordonnée par coordonnée. Pas de synchronisation avec l'hôte,
        pas de décomposition supplémentaire. Les statistiques ne sont calculées que lorsqu'on les lit.
  \item \textbf{Ce qui est comparé dans E16.} Les trois règles (7 valeurs de $q$ chacune), la famille A
        ($\max(s^\star, \lambda)$, 5 valeurs de $\lambda$) et le correctif E14 ($s^\star + \lambda$,
        les mêmes 5 valeurs, relancé dans E16), sur \fic{cnn_gn_cifar}, \fic{vit_micro_cifar} et
        \fic{cct_2_3x2_cifar}, avec \fic{ekfac} et \fic{tekfac} et 5 graines. Les règles de décision
        sont pré-enregistrées dans \fic{docs/reports/plan_lambda_dominance.md}, section E16.
\end{itemize}

\end{document}
